\subsection{ELEMENTARY FILTERS}

DEFINITION 7. Let \((x_{n})_{n\in \mathbb{N}}\) be an infinite sequence of elements of a set \(\mathbf{X}\) . The elementary filter associated with the sequence \((x_{n})\) is the filter generated by the image of the Fréchet filter (no. 1) by the mapping \(n\to x_n\) of \(\mathbf{N}\) into \(\mathbf{X}\) . It comes to the same thing to say that the elementary filter associated with \((x_{n})\) is the set of subsets \(\mathbf{M}\) of \(\mathbf{X}\) such that \(x_{n}\in \mathbf{M}\) except for a finite number of values of \(n\) . If \(\mathbf{S}_n\) denotes the set of all \(x_{p}\) such that \(p\geqslant n\) , then the sets \(\mathbf{S}_n\) form a base of the elementary filter associated with the sequence \((x_{n})\) .

The elementary filter associated with an infinite subsequence of a sequence \((x_{n})\) is finer than the elementary filter associated with \((x_{n})\) (cf.~Exercise 15).

By definition, every elementary filter has a countable base. Conversely: PROPOSITION 11. If a filter F has a countable base, it is the intersection of the elementary filters which are finer than F.

Let us arrange the countable base of \(\mathfrak{F}\) as a sequence \((\mathbf{A}_n)_{n\in \mathbb{N}}\) ; if we put

\[
\mathrm{B} _ {n} = \bigcap_ {p = 0} ^ {n} \mathrm{A} _ {p},
\]

then the \(B_{n}\) again form a base of F (no. 3, Proposition 3) and we have \(B_{n+1} \subset B_{n}\) for each n.~Let \(a_{n}\) be any element of \(B_{n}\) for each \(n \in N\) ; then it is clear that F is coarser than the elementary filter associated with \((a_{n})\) . Hence the intersection S of the elementary filters which are finer than F exists and is finer than F; if S is strictly finer than F there exists a set \(M \in S\) such that \(B_{n} \cap CM \neq \varnothing\) for each n; if \(b_{n} \in B_{n} \cap CM\) , the elementary filter associated with the sequence \((b_{n})\) is finer than F and does not contain M. This contradicts the definition of S.

Remark. A filter which is coarser than a filter with a countable base need not possess a countable base; for example, if X is an uncountable infinite set, then the filter consisting of the complements of finite subsets of X has no countable base (otherwise the set of finite subsets of X would be countable, contrary to assumption); nevertheless this filter is coarser than every elementary filter associated with an infinite sequence of distinct elements of X.

